\documentclass[10pt,a4paper]{amsart}
\usepackage[margin=19mm]{geometry}
\usepackage{amsmath,amssymb,mathtools}
\usepackage[scr=rsfs,cal=euler]{mathalfa}
\usepackage{xfrac}
\usepackage[unicode,hidelinks,bookmarks=false]{hyperref}
\newcommand{\ie}{i.e.\ }
\newcommand{\cf}{cf.\ }
\newcommand{\resp}{respectively\ }
\newcommand{\Subsection}{\subsection}
\providecommand{\nobreakdash}{\nobreak\hbox{-}}
\setlength{\emergencystretch}{2em}
\multlinegap0pt
\usepackage{xcolor}
\usepackage{amssymb}
\usepackage{bm}
\usepackage{mathtools}
\newtheorem{proposition}{Proposition}
\title{Neural Preconditioned Born Series: A Metric-Matched Framework for Learning-based Preconditioners}
\author{Juntao Wang, Jiwei Jia, Xinliang Liu}
\date{Source: arXiv:2603.18527v4 (CC BY 4.0)}
\begin{document}
\maketitle

\noindent\emph{Source and licence.} Neural Preconditioned Born Series: A Metric-Matched Framework for Learning-based Preconditioners. Version arXiv:2603.18527v4 (CC BY 4.0), distributed by arXiv under the Creative Commons Attribution 4.0 International licence, http://creativecommons.org/licenses/by/4.0/, as recorded in the arXiv licence metadata for this exact version and shown on the abstract page https://arxiv.org/abs/2603.18527v4. Full licence notice: \texttt{licenses/arxiv-2603.18527-CC-BY-4.0.txt}.\par\smallskip

\noindent\textit{Editorial scope.} Counted unit: the article's proposition prop:fft\_all\_cases (source0.tex lines 941--944). The article's complete original proof of it is delivered with it (source0.tex lines 946--958, one \texttt{proof} environment of 13 lines). No mathematics is written, repaired or re-derived: the statement and the proof are the article's own lines, in the article's own notation, and no retained line is substituted or re-flowed.  No further source-local context is needed: every cross-reference of every retained line resolves inside the two delivered spans.  Nothing was omitted from any retained span, so the package carries no omission record.  No retained line contains a citation, so the wrapper carries no bibliography and no bibliography entry.  No retained line contains a figure environment, an image-inclusion command or drawing code, so no image is delivered and the package declares no asset dependency.  Editorial formatting only, in addition: an AMS \texttt{amsart} preamble that restates the article's own declarations and macros used by the retained lines, namely the environments \texttt{proposition} on separate counters, together with the standard packages those lines need.  The retained environments print in this document as Proposition 1 in order of appearance, whereas the article prints the same items with the numbering of its own full document.  The complete native source of the article as distributed in the arXiv e-print for this exact version is bundled unchanged as \texttt{sources/196701/source0.tex}.
\begin{proposition}[Transform-diagonalizable Green operators]
\label{prop:fft_all_cases}
On a uniform grid, if the chosen reference operator has constant coefficients and compatible boundary conditions (periodic, or transform-compatible Dirichlet/Neumann), then its discrete matrix is diagonalizable by a fast trigonometric transform. Consequently, applying the corresponding discrete Green operator costs $\mathcal{O}(N\log N)$.
\end{proposition}

\begin{proof}[Constructive verification]
Let $L_0$ be any constant-coefficient reference operator and $G_0=L_0^{-1}$. Under periodic boundary conditions, the discrete $L_0$ is block-circulant with circulant blocks, hence
\[
L_0 = F^{-1}\Lambda F,\qquad
G_0 = F^{-1}\Lambda^{-1}F,
\]
where $F$ is the multidimensional DFT matrix and $\Lambda$ contains the symbol values. Therefore
\[
G_0 q = \mathcal{F}^{-1}\!\left(\frac{\widehat{q}}{\lambda(\xi)}\right),
\]
implemented by FFT/IFFT in $\mathcal{O}(N\log N)$.
For Dirichlet or Neumann boundaries, the same statement holds with sine/cosine transforms (DST/DCT), which are FFT-based and have the same complexity order.
\end{proof}

\end{document}
